\chapter{Homeostasis and functional balance}\label{ch:homeostasis}

\section{The tank level stays steady while water moves}
Stand beside the clinic's tank for a few minutes. Water enters through one pipe and leaves through another. The level can remain nearly unchanged even though the water present is continually being replaced. A photograph would show stability. A flow meter would show activity. Both observations can be correct.

Now close the inlet while leaving the outlet open. The level begins to fall. A caretaker notices and restores the inlet. If the available inflow can match the withdrawal, the level may recover. If the source is empty or the pump cannot provide enough, noticing the problem will not be sufficient. A useful account of regulation must include both information and physical capability.

This simple story introduces homeostasis: the maintenance of selected conditions compatible with continuing function despite ongoing processes and disturbances. The word does not require complete immobility, perfect constancy or a universal target shared by every system. It directs attention to what is maintained, how deviations are detected and which processes respond.

In physiology, McFarland and colleagues developed and validated a teaching framework for homeostasis through negative feedback at the organism level. Their framework concerns regulated variables and the mechanisms that influence them. It is a conceptual educational source, not a quantitative model of society or proof that every ecosystem follows one feedback design.\cite{intro-homeostasis}

The analogy with a resource network is useful only if we preserve that difference. A body and a town are not the same kind of system. A town contains people with rights, conflicting purposes and institutions. We borrow a question about continuing function; we do not infer a social controller from biological vocabulary.

\section{What is being maintained?}
Before calling something homeostatic, identify the variable and the function. In the tank example, the stock supports a declared service. In a different example, the relevant condition might be temperature, pressure or an equipment state. These quantities have different units and mechanisms. They cannot be merged by describing all of them as ``balance''.

The variable may also be only a proxy. A measured tank level can indicate available water quantity, but not necessarily water quality. A measured room temperature may not capture every condition relevant to storage. If a proxy is used, the relationship to function needs evidence. Good regulation of the proxy can coexist with poor regulation of what was actually intended.

For a model, we therefore distinguish the represented condition from the wider purpose. The state variable is what the mathematics contains. The functional interpretation explains why that variable was chosen. The evidence for the interpretation belongs to the application. Keeping those layers visible allows a model to be corrected when its proxy proves inadequate.

The same care applies to a reference. A declared reference stock may be convenient for a teaching account, but an applied reference needs justification. It might depend on the intended service, physical design or an agreed planning criterion. The reference is not automatically a natural constant, and different places may appropriately have different values.

\section{Negative feedback in ordinary language}
In a negative-feedback arrangement, a detected departure elicits an influence that tends to oppose that departure under the stated conditions. If the tank falls, the arrangement increases inflow or reduces outflow. If it rises beyond the intended level, the arrangement decreases inflow or permits additional outflow. The word ``negative'' concerns opposition to a departure; it does not mean morally bad.

This description contains several assumptions. The departure must be observed with sufficient timeliness. The relevant action must be available. The sign of the response must be appropriate. The response must not be so delayed or strong that it creates a different problem. A diagram showing an arrow back to a comparator does not prove that all these conditions hold.

Consider a hypothetical caretaker who responds to yesterday's low reading by opening the valve fully today. If the tank has already filled through another route, the delayed response may overshoot. A different caretaker may react immediately but lack enough source stock. The first arrangement has an information problem; the second has an action-capacity problem. Both can fail despite an intention to restore the same condition.

This book does not derive a controller that resolves such dynamic questions. It develops a potential and exact finite action account. The account may help describe whether a particular accepted action moves the represented state toward or away from its reference. A theorem about that description is different from a theorem that repeated decisions produce stable regulation.

\begin{cautionbox}{A restoring description is not a restoring mechanism}
Drawing a landscape with its lowest point at a reference does not make the physical system travel downhill. Movement requires an action law, available routes, permission and resources. A proof about the landscape alone establishes no recovery rate or long-run homeostatic performance.
\end{cautionbox}

\section{Functional balance is not equality everywhere}
Ridge and Vale need not hold identical stocks in every useful arrangement. A storage facility, a clinic and a transit point may perform different functions. Even in a single-resource model, the declared reference may assign different stocks to different nodes. Equal numbers can be a convenient teaching choice without being a universal law.

Suppose a hypothetical two-place system contains twenty units and uses a reference $(14,6)$. That reference is compatible with the total. A state $(10,10)$ has equal stocks but does not match the declared reference. The physical equality of the two numbers therefore cannot settle functional balance. We must ask why the reference was chosen and whether it suits the intended application.

The scales matter as well. A difference of two units may have different significance at two places. The later Gaussian construction divides each deviation by a declared positive scale before squaring it. This creates a dimensionless comparison within the model. It does not prove that the scale is objectively correct or that unlike services have become interchangeable.

The phrase ``functional balance'' will therefore refer to an arrangement's relation to its declared functional description, rather than to an insistence that every coordinate be equal. Where the description is incomplete, the conclusion must remain incomplete. The mathematical clarity lies in making the declaration inspectable.

\section{A stable stock is not necessarily a healthy arrangement}
Imagine a tank sealed shut with ten units inside. Its level is constant. If the intended function is to supply the clinic, the absence of delivery may be a failure. Now imagine a tank whose level stays at ten because inflow matches useful withdrawal. The same stock observation accompanies a different process.

This distinction is central to any future study of sustained service. A rule that does nothing can leave an action-only world unchanged. If no disturbance or consumption exists, the world may remain within its initial conditions indefinitely. That is not evidence that the rule can meet changing needs. The test must contain the processes relevant to the claim.

Conversely, a temporarily changing stock is not automatically a failure. A deliberate transfer can reduce a source stock while improving the represented arrangement. A service can require withdrawal. Maintenance can temporarily interrupt operation. The judgement depends on the declared purpose and constraints, rather than on a general preference for every variable to remain constant.

Homeostasis is therefore better understood through the maintenance of function over an appropriate interval than through the absence of visible change. The interval and the function must both be stated. A snapshot can support a state description; it cannot by itself establish recovery after disturbance or continuing provision.

\section{A physical balance before a valuation}
For the tank, an elementary stock account is
\begin{equation}
\text{new stock}=\text{old stock}+\text{inflow}-\text{outflow},
\end{equation}
where every term uses the same physical unit and the flows refer to the declared interval. This is a balance statement. It does not yet say whether the new stock is desirable.

If the old stock is ten units, inflow is four and outflow is four, the new stock is ten. If inflow is two and outflow is four, the new stock is eight. The arithmetic is simple, but its interpretation depends on the boundary. If some water remains in a connecting pipe, the account needs to say whether that water is included in inflow, transit or neither yet.

The fixed-total teaching model later in the book considers internal transfers among represented places. A lossless transfer subtracts a quantity from one place and adds the same quantity to another. External inflow and withdrawal are not silently present. When discussed, they are declared separately, because they change the physical accounting conditions.

This ordering is deliberate. We first establish what changed. We then evaluate that change through a declared potential. If we reverse the order, a desirable interpretation can tempt us to overlook a missing physical term. No reduction in represented burden excuses an unexplained creation of resource.

\section{The role of reserves and limits without importing a controller}
Ordinary life contains limits. A tank cannot contain a negative quantity, a pipe has physical capabilities, and essential functions can require availability. Recognising these facts does not require importing every historical threshold model into this introduction. The present edition develops a smooth reference-based potential and keeps physical permission as a separate declaration.

A future application may impose a source reserve or a storage-capacity limit. Such a restriction would constrain the allowed action set. It would not automatically become a term in the potential, and its numerical value could not be inferred merely from a desire for safety. A reserve protects a specified purpose under specified assumptions; its meaning must be stated.

This distinction helps prevent one of the most common modelling confusions. A permission rule that limits export does not necessarily direct resource toward a place that needs it. Refusing to empty a source and restoring a destination are different actions. The former may be necessary for safety while remaining insufficient for service.

The Gaussian landscape likewise should not be read as a complete permission rule. An action can move toward the reference while violating a physical restriction. Another can be physically permitted while moving away from the reference. The account becomes clearer when permission and valuation are checked separately.

\section{Why a smooth potential is a useful teaching choice}
The later potential assigns burden continuously as a state moves away from its declared reference. Its quadratic form allows exact finite arithmetic and makes curvature visible. A derivative describes the local slope, while a finite difference describes a completed finite change. The relation between these objects can be derived rather than guessed.

This is a mathematical advantage, not a claim that every biological or social response is quadratic. A convenient model can be useful for establishing concepts even when a future application needs a different representation. The correct question is what the construction assumes and what follows from those assumptions.

It also gives us a clean way to discuss overshooting. Moving toward a reference can reduce burden up to a point; continuing past it can eventually undo the improvement. A positive initial slope along an action direction therefore does not justify multiplying that slope by an arbitrary finite quantity. Later chapters calculate the exact correction.

At this stage, the essential insight is qualitative. ``More action'' is not the same as ``more restoration''. The state changes while the action proceeds, so the value of an additional increment changes as well. A finite account must evaluate the actual change rather than freeze the first impression indefinitely.

\section{Several functions can make competing demands}
Suppose the clinic and the workshop draw from the same store. The clinic needs water for current service; the workshop needs water to repair the equipment that will support later service. Sending everything to one destination may impair the other function. Calling both needs homeostatic does not remove the competition.

A physical model can expose the shared constraint: the same source quantity cannot be sent twice. A potential can express a declared evaluation of the resulting states. A permission rule can exclude physically invalid proposals. Yet a choice among valid competing uses still depends on the declared policy and purpose. No general biological metaphor supplies that choice automatically.

The example also shows why independent local proposals require a common acceptance step when they contend for the same stock. Each proposer may see an available quantity and request it. Accepting all requests separately could exceed the source. The final account must evaluate the vector actually accepted, because the physical world receives the combined action rather than each proposal's private imagined world.

Later chapters will examine this shared-field issue mathematically. Here it is enough to see its functional meaning. Maintaining one variable cannot be treated as success if the action quietly disables another declared essential function. Whether the representation captures both functions is an application question; whether the resource is double-spent is an accounting question. Both deserve a clear answer.

\section{The caretaker's questions}
Before using a homeostatic description, our imaginary caretaker can ask a sequence of practical questions. Which function are we trying to keep available? Which state variable represents a condition relevant to it? What reference is being used? Which observations support the current state? Which actions are permitted? What will be measured after the action?

These questions do not demand that the caretaker solve every philosophical issue before opening a valve. They demand that a research model state the assumptions on which its conclusion depends. A routine operational rule can be simple because much of that design work has already been done. A new theory must expose that work rather than borrow the appearance of simplicity.

The caretaker should also ask what happens when the assumptions fail. If a sensor is unavailable, the account needs a declared response. If the resource changes quality, the homogeneous model may no longer apply. If a new service is added, the old reference may need review. Such review changes the model prospectively; it should not rewrite past actions under an undisclosed new evaluation.

\section{Guided problems and answers}
\subsection{The unchanging tank}
Two tanks remain at ten units for an hour. One is sealed. The other supplies four units while receiving four. Do the equal final stocks establish equal service?

No. The state snapshots agree, but the flow histories differ. A service claim requires a record of the relevant delivery. The sealed tank illustrates why preservation of a stock alone is an insufficient success criterion for a system intended to provide something.

\subsection{The well-drawn landscape}
A model has a potential with a unique minimum at its declared reference. Does a state away from the reference necessarily move toward it?

No. The potential is an evaluation. A state changes only through the declared physical dynamics or actions. Available routes, restrictions and a selection rule determine whether and how movement occurs. A unique minimum does not imply a particular trajectory.

\subsection{The export limit}
A rule prevents Ridge from sending more than a declared quantity. Does that guarantee that Vale receives enough?

No. The rule limits an action at the source. It does not by itself select a destination, create a route or require delivery. Service requires additional conditions. This is why the language of protection should always identify what is protected and by which mechanism.
